\scanpage{306}
\mspage{256}

\uline{Cor. 1} Considérons $p\colon\Psi\mathrel{\unsure{-|\to}}e$, et le foncteur
\[
p_\Psi^*\colon\mathrm{Ho}_W(e)=\mathrm{Ho}_W\longrightarrow\mathrm{Ho}_W(\Psi),
\]
déduit par passage aux cat. de fractions du foncteur
\diagram{\(\displaystyle\begin{aligned}
M&\longrightarrow M(\Psi)\\
X&\longmapsto\left(\begin{tikzcd}[ampersand replacement=\&,row sep=small,column sep=small]
X\arrow[d]\&\\\deleted{S}X\&X\arrow[l]
\end{tikzcd}\right).
\end{aligned}\)}

Ce foncteur admet \deleted{\unsure{du}} le \deleted{adjoint \unsure{à}} foncteur $\varepsilon_3^*\unsure{i|e}_*$ précédent comme adjoint à droite, i.e.
\[
p_{\Psi*}=\varepsilon_3^*i_*\colon\mathrm{Ho}_W(\Psi)\longrightarrow\mathrm{Ho}_W.
\]

\uline{Dém} On peut regarder $p_\Psi$ comme le composé
\[
\Psi\mathrel{\overset{i}{\unsure{-|\longrightarrow}}}\Psi^{\sim}\xrightarrow{p_{\widetilde\Psi}}e,
\]
et on a donc
\[
p_\Psi^*\simeq i^*(p_{\widetilde\Psi})^*.
\]

Donc l'existence et la détermination d'un adjoint à droite pour $p_\Psi^*$ \unsure{résultent de|résulteront de} \deleted{\unsure{de}}l'existence et de la détermination d'adj. à droite $i_*,(p_{\widetilde\Psi})_*$ de $i^*,(p_{\widetilde\Psi})^*$, via
\[
(p_\Psi)_*\simeq(p_{\widetilde\Psi})_*i_*.
\]
% archive sheet number: 306
